% The quantum CNOT gate on the four basis states.
% Source: book/tex/quantum/ (Qcircuit diagram, ported to quantikz).
\documentclass[border=2mm]{standalone}
\usepackage{tikz}
\usepackage{amssymb}
\usepackage{quantikz}
\providecommand{\ket}[1]{\left\lvert #1\right\rangle}
\providecommand{\xup}{\ket\rightarrow}
\providecommand{\xdown}{\ket\leftarrow}
\begin{document}
\begin{quantikz}
\lstick{$\ket0$} & \ctrl{1} & \rstick{$\ket0$}\qw \\
\lstick{$\ket0$} & \targ{} & \rstick{$\ket0$}\qw
\end{quantikz}
\hspace{6em}
\begin{quantikz}
\lstick{$\ket0$} & \ctrl{1} & \rstick{$\ket0$}\qw \\
\lstick{$\ket1$} & \targ{} & \rstick{$\ket1$}\qw
\end{quantikz}
\hspace{6em}
\begin{quantikz}
\lstick{$\ket1$} & \ctrl{1} & \rstick{$\ket1$}\qw \\
\lstick{$\ket0$} & \targ{} & \rstick{$\ket1$}\qw
\end{quantikz}
\hspace{6em}
\begin{quantikz}
\lstick{$\ket1$} & \ctrl{1} & \rstick{$\ket1$}\qw \\
\lstick{$\ket1$} & \targ{} & \rstick{$\ket0$}\qw
\end{quantikz}

\end{document}
